\begin{description}
\item[(a)] Step 1: Determine the GST at the time of observation.

  Let us denote the GST at 00:00 UT on 1st January 2020 as
  $\mathrm{GST}_0$. Total 268 days have elapsed since start of the year till
  $0^{\mathrm{h}}$ UT on 25 September. \hfill(1.0)

  At 21:00 for GMT+8 timezone, UT will be
  $21^{\mathrm{h}} - 8^{\mathrm{h}} = 13^{\mathrm{h}}$. Hence,
  \begin{align*}
    GST &= GST_0 + \Delta t\\
        &= 6^{\mathrm{h}}40^{\mathrm{m}}30^{\mathrm{s}}
           + \frac{268 \times 24}{365.2422} + \frac{13 \times 24}{23.9344}\\
        &\approx 6^{\mathrm{h}}40^{\mathrm{m}}30^{\mathrm{s}}
           + 17^{\mathrm{h}}36^{\mathrm{m}}36^{\mathrm{s}}
           + 13^{\mathrm{h}}2^{\mathrm{m}}8^{\mathrm{s}}\\
        &= 37^{\mathrm{h}}19^{\mathrm{m}}14^{\mathrm{s}}
         = \mathbf{13^{h}19^{m}14^{s}}
  \end{align*}
  \hfill(3.0)

  As longitude of observer is $120.4^\circ$, the LST of observation is,
  \begin{align*}
    LST &= GST + 24^{\mathrm{h}} \times \frac{120.4^\circ}{360^\circ}\\
        &= 13^{\mathrm{h}}19^{\mathrm{m}}14^{\mathrm{s}}
           + 8^{\mathrm{h}}1^{\mathrm{m}}36^{\mathrm{s}}\\
        &= \mathbf{21^{h}20^{m}50^{s}}
  \end{align*}
  \hfill(2.0)

\item[(b)] For horizontal coordinates

  % FIGURE NEEDED: spherical-triangle sketch with the zenith, the radiant
  % (A,a), and the initial (A_i,a_i) and final (A'_i,a'_i) points of the i-th
  % meteor, with the angle theta at the meteor track; 2020 (GeCAA) theory
  % solutions, page 18 (problem 7).
  \hfill(1.0)

  The spherical triangle has the radiant, initial position and final position
  of the $i$-th meteor to be $(A, a), (A_i, a_i)$ and $(A'_i, a'_i)$, where $a$
  and $A$ denotes the altitude and azimuth respectively.

  Using the four-parts (cotangent) equation on the left and right triangles, we
  have, for both meteors:
  \begin{align*}
    \cos(90 - a_i)\cos(A'_i - A_i)
      &= \sin(90 - a_i)\cot(90 - a'_i) - \sin(A'_i - A_i)\cot\theta\\
    \therefore \sin a_i \cos(A'_i - A_i)
      &= \cos a_i \tan a'_i - \sin(A'_i - A_i)\cot\theta
  \end{align*}
  Similarly,
  \begin{align*}
    \cos(90 - a_i)\cos(A_i - A)
      &= \sin(90 - a_i)\cot(90 - a) - \sin(A_i - A)\cot(180 - \theta)\\
    \therefore \sin a_i \cos(A_i - A)
      &= \cos a_i \tan a + \sin(A_i - A)\cot\theta
  \end{align*}
  \hfill(2.0)

  We can then eliminate $\cot\theta$ to yield:
  \[
    \frac{\sin a_i \cos(A_i - A)}{\sin(A_i - A)}
      - \frac{\cos a_i \tan a}{\sin(A_i - A)}
    = \frac{\cos a_i \tan a'_i}{\sin(A'_i - A_i)}
      - \frac{\sin a_i \cos(A'_i - A_i)}{\sin(A'_i - A_i)}
    = \cot\theta
  \]
  multiplying the whole equation by
  $\dfrac{\sin(A'_i - A_i)\sin(A_i - A)}{\cos a_i}$
  \begin{align*}
    \tan a_i \cos(A_i - A)\sin(A'_i - A)
      - \tan a \sin(A'_i - A)
      % sic: the printed solution has sin(A'_i - A) in the two left-hand
      % terms; mathematically it should be sin(A'_i - A_i).
      &= \tan a'_i \sin(A_i - A)
       - \tan a_i \cos(A'_i - A_i)\sin(A_i - A)\\
    \tan a_i \sin[(A'_i - A_i) + (A_i - A)]
      &= \tan a'_i \sin(A_i - A) + \tan a \sin(A'_i - A_i)\\
    \tan a_i \sin(A'_i - A)
      &= \tan a'_i \sin(A_i - A) + \tan a \sin(A'_i - A_i)
  \end{align*}
  \hfill(3.0)

  Plugging in $i = 1, 2$, and eliminating $\tan a$ by division, we have
  \[
    \frac{\tan a_1 \sin(A'_1 - A) - \tan a'_1 \sin(A_1 - A)}
         {\tan a_2 \sin(A'_2 - A) - \tan a'_2 \sin(A_2 - A)}
    = \frac{\sin(A'_1 - A_1)}{\sin(A'_2 - A_2)}
    \overset{\text{def}}{=} k
  \]
  where $k$ is the RHS defined above (which is a known value). Expanding the
  equation above and dividing by $\cos A$, gathering terms of $\tan A$, we get
  \[
    \tan A = \frac{\tan a_1 \sin A'_1 - \tan a'_1 \sin A_1
                   - k\tan a_2 \sin A'_2 + k\tan a'_2 \sin A_2}
                  {\tan a_1 \cos A'_1 - \tan a'_1 \cos A_1
                   - k\tan a_2 \cos A'_2 + k\tan a'_2 \cos A_2}
  \]
  \hfill(2.0)

  and
  \[
    \tan a = \frac{\tan a_1 \sin(A'_1 - A) - \tan a'_1 \sin(A_1 - A)}
                  {\sin(A'_1 - A_1)}
  \]

  Plugging in the respective values from the question,
  \begin{align*}
    A_1 &= 0^\circ  & a_1 &= 15^\circ\\
    A'_1 &= 90^\circ & a'_1 &= 0^\circ\\
    A_2 &= 210^\circ & a_2 &= 23.5^\circ\\
    A'_2 &= 255^\circ & a'_2 &= 75^\circ
  \end{align*}
  We get
  \begin{align*}
    k &= \frac{\sin(A'_1 - A_1)}{\sin(A'_2 - A_2)}
       = \frac{\sin 90^\circ}{\sin 15^\circ}\\
    k &= 1.414\\
    \tan a_1 \sin A'_1 &= 0.268\\
    \tan a'_1 \sin A_1 &= 0\\
    k\tan a_2 \sin A'_2 &= -0.594\\
    k\tan a'_2 \sin A_2 &= -2.639\\
    \tan a_1 \cos A'_1 &= 0\\
    \tan a'_1 \cos A_1 &= 0\\
    k\tan a_2 \cos A'_2 &= -0.159\\
    k\tan a'_2 \cos A_2 &= -4.571\\
    \therefore \tan A &= \frac{0.268 - 0 + 0.594 - 2.639}
                              {0 - 0 + 0.159 - 4.571} = 0.403\\
    \therefore A &= 21.94^\circ\\
    \text{OR } A &= 201.94^\circ
  \end{align*}
  \hfill(4.0)

  However, the student should realize that for $A = 21.94^\circ$, the radiant
  is in the path of the first meteor, which is unrealistic. Hence, it is more
  likely that the azimuth takes the second value. \hfill(1.0)

  Plugging in for the altitude, we find the coordinate of the radiant to be
  \begin{align*}
    \tan a_1 \sin(A'_1 - A) &= -0.249\\
    \tan a'_1 \sin(A_1 - A) &= 0\\
    \sin(A'_1 - A_1) &= 1\\
    \therefore \tan a &= (-0.2485 - 0)/1 = -0.2485\\
    a &= -13.96^\circ
  \end{align*}
  \hfill(2.0)

  As the meteor track origin and end points are rounded to integer degrees,
  the horizontal coordinates of the radiant will be
  \[
    (\mathbf{A}, \mathbf{a}) = (202^\circ, -14^\circ)
  \]
  \hfill(1.0)

  \textbf{Alternative solution}

  In the vector representation, let the starting and final locations of the
  meteroids be $\vec{u}_i$ and $\vec{v}_i$ respectively. For the subsequent
  calculations, the vectors are most conveniently represented in Cartesian
  coordinates. Letting the $x$-axis point east, $y$-axis north, and $z$-axis
  towards the zenith, we have
  \[
    \vec{u} = (u_x, u_y, u_z) = (\cos a \sin A, \cos a \cos A, \sin a)
    \tag{1}
  \]
  for a generic vector $\vec{u}$ with azimuth $A$ and altitude $a$. We can
  therefore compute
  \begin{align*}
    \vec{u}_1 &= (\cos 15^\circ \sin 0^\circ,
                  \cos 15^\circ \cos 0^\circ,
                  \sin 15^\circ) = (0, 0.9659, 0.2588)\,,\\
    \vec{u}_2 &= (\cos 23.5^\circ \sin 210^\circ,
                  \cos 23.5^\circ \cos 210^\circ,
                  \sin 23.5^\circ) = (-0.4585, -0.7942, 0.3988)\,,\\
    \vec{v}_1 &= (\cos 0^\circ \sin 90^\circ,
                  \cos 0^\circ \cos 90^\circ,
                  \sin 0^\circ) = (1, 0, 0)\,,\\
    \vec{v}_2 &= (\cos 75^\circ \sin 255^\circ,
                  \cos 75^\circ \cos 255^\circ,
                  \sin 75^\circ) = (-0.25, -0.0670, 0.9659)\,.
  \end{align*}
  \hfill(3.0)

  The meteors move along great arcs, and the two intersections of the great
  arcs correspond to the two possible locations of the radiant. A great arc is
  uniquely defined by its normal vector, denote it with $\vec{n}_i$. The
  normal vector is conveniently calculated through the cross product of the
  starting and final position, $\vec{n}_i = \vec{u}_i \times \vec{v}_i$. The
  cross product can be calculated via the determinant of the following matrix
  \[
    \vec{u}_i \times \vec{v}_i =
    \begin{vmatrix}
      i & j & k\\
      u_{ix} & u_{iy} & u_{iz}\\
      v_{ix} & v_{iy} & v_{iz}
    \end{vmatrix}
    = (u_{iy}v_{iz} - u_{iz}v_{iy},\,
       u_{iz}v_{ix} - u_{ix}v_{iz},\,
       u_{ix}v_{iy} - u_{iy}v_{ix}).
  \]
  \hfill(4.0)

  This yields
  \begin{align*}
    \vec{n}_1 &= (0, 0.2588, -0.9659)\,,\\
    \vec{n}_2 &= (-0.7404, 0.3432, -0.1678)\,.
  \end{align*}
  \hfill(1.0)

  Every point on the great circle $i$ is perpendicular to $\vec{n}_i$. This
  means that the intersections of the two great circles are perpendicular to
  both $\vec{n}_1$ and $\vec{n}_2$. The vector parallel to the intersections
  can therefore be calculated by the cross product of $\vec{n}_1$ and
  $\vec{n}_2$. Let the vector corresponding to the intersection be $\vec{m}$.
  Then
  \begin{align*}
    \vec{m} &= \vec{n}_1 \times \vec{n}_2 & \text{(3.0)}\\
            &= (0.3626, 0.9002, 0.2412)\,. & \text{(1.0)}
  \end{align*}

  Now all that's left is to calculate the azimuth and altitude corresponding
  to $\vec{m}$. Note that $-\vec{m}$ is also an intersection. The equations
  for the altitude and azimuth can be reverse-engineered from equation (1):
  \[
    a = \arctan\left(\frac{z}{\sqrt{x^2 + y^2}}\right), \qquad
    A = \arctan\left(\frac{x}{y}\right).
  \]
  \hfill(1.0)

  One has to take care that the arctans are signed. This gives the final
  values of the intersection points to be
  \[
    (A, a) = (22^\circ, 14^\circ)
    \quad\text{and}\quad
    (A, a) = (202^\circ, -14^\circ).
  \]
  \hfill(2.0)

  As in the first solution, the $A = 22^\circ$ solution can be eliminated due
  to it being on the path of the first meteor. This gives the final answer of
  \[
    (\mathbf{A}, \mathbf{a}) = (202^\circ, -14^\circ)
  \]
  \hfill(1.0)

\item[(c)] Now, it is easy to convert the horizontal coordinates of the
  radiant to equatorial coordinates.

  % FIGURE NEEDED: celestial-sphere diagram with zenith Z, pole P, points
  % N, E, S, W, O, the radiant J, and the arcs (90-phi), (90-delta), (90-a),
  % 360-A and hour angle H; 2020 (GeCAA) theory solutions, page 21
  % (problem 7).
  \hfill(1.0)

  Using the cosine rule:
  \begin{align*}
    \sin\delta &= \sin\phi \sin a + \cos\phi \cos a \cos A\\
      &= \sin(23.5^\circ)\sin(-13.4^\circ)
         + \cos(23.5^\circ)\cos(-13.96^\circ)\cos(201.94^\circ)\\
      % sic: the official solution writes -13.4 deg in two places here and
      % below where the altitude is a = -13.96 deg.
      &= -0.9217\\
    \therefore \delta &= -67.2^\circ
  \end{align*}
  \hfill(2.0)
  \begin{align*}
    \cos H &= \frac{\sin a - \sin\delta \sin\phi}{\cos\phi \cos\delta}\\
      &= \frac{\sin(-13.4^\circ) - \sin(-67.2^\circ)\sin(23.5^\circ)}
              {\cos(23.5^\circ)\cos(-67.2^\circ)} = 0.3551\\
    \therefore H &= 4^{\mathrm{h}}36^{\mathrm{m}}47^{\mathrm{s}}
  \end{align*}
  \hfill(1.5)

  From the relationship between the LST and the Hour Angle
  ($H = LST - \alpha$), we get the equatorial coordinates of the radiant to be
  \begin{align*}
    \alpha &= 21^{\mathrm{h}}20^{\mathrm{m}}50^{\mathrm{s}}
              - 4^{\mathrm{h}}36^{\mathrm{m}}47^{\mathrm{s}}\\
           &= 16^{\mathrm{h}}44^{\mathrm{m}}3^{\mathrm{s}}
  \end{align*}
  \hfill(1.5)

  Thus, the equatorial coordinates of the radiant are,
  \[
    (\alpha, \delta) = (16^{\mathrm{h}}44^{\mathrm{m}}, -67.2^\circ)
  \]

\item[(d)] Triangulum Australes \hfill(2.0)
\end{description}
